\ctexset{chapter/number={\Roman{chapter}}}
\chapter{土壤}
\subsection{土壤成分}
土壤是疏松散碎和多孔的物质，是人类从事农业生产的一项重要资源。
它是由固体、液体和气体三部分组成的复杂体系。
土壤里的固体包括多种化学组成不同的矿物质和有机质；土壤中的液体主要是土壤水分，其中溶解着多种盐类；主壤中的气体就是空气。
水和空气常常共存于土壤孔隙中，水多气就少，气多水就少。
土壤中还生活着多种微生物。

土壤里的矿物质是由岩石、矿物经过长期的风吹、日晒、雨淋、生物生长等过程逐渐形成的，其中含有铝、铁、钙、镁、钾、硅、磷等多种元素，组成铝硅酸盐和氧化物等复杂化合物。

土壤里的有机质比矿物质少得多。
土壤有机质主要来自动植物残体、施入的有机肥料以及腐殖质和微生物等。
土壤有机质主要是由碳、氢、氧、氮以及磷、硫等元素组成的。
土壤有机质在土壤水分、空气、动物以及微生物的作用下，除小部分转变为腐殖质外，大部分由复杂的有机化合物变成简单的无机化合物，如水、二氧化碳、氨等。
在这一过程中，有机质经分解可以释放出植物能够吸收的养分。
土壤有机质在提高土壤肥力中起着重要的作用。

\subsection{土壤胶体}
组成土壤胶体的物质有无机物也有有机物。
土壤里的无机胶体主要是铝硅酸盐及铁、铝的含水氧化物等。
土壤里的有机胶体主要是土壤腐殖质。
实际上，这两种胶体常常结合在一起。
研究土壤胶体对于揭露土壤中化学现象的本质有极为重要的意义。

\subsubsection{土壤胶体的带电性质}
土壤胶体的胶体微粒具有很大的表面积，能吸附离子，因而带有电荷。
实践证明，土壤胶体主要带负电荷。

\subsubsection{土壤胶体的离子吸附和离子交换过程}
土壤胶体既带负电荷，又有很大的表面积，因而具有吸附周围液体中的阳离子的作用。
这些被吸附的阳离子能够跟土褰溶液里的阳离子起交换作用。
例如，施入酸性土壤里的石灰，不但可以中和土壤溶液里的酸性，而且可以交换土壤胶体吸附着的氢离子和铝离子。
\begin{align*}
\text{（土壤胶体）}\!
\begin{array}{c} \ce{H+}\\[-5pt] \ce{H+} 
\end{array} 
\ce{ + Ca^2+ + 2OH- }&\ce{<=>}\text{（土壤胶体）}\!\ce{Ca^2+ + 2H2O} \\
\text{（土壤胶体）}\!
\begin{array}{c} \ce{Al^3+}\\[-5pt] \ce{Al^3+} 
\end{array} 
\ce{ + 3Ca^2+ + 6OH- }&\ce{<=>}\text{（土壤胶体）}\!
\begin{array}{c} \ce{Ca+}\\[-5pt] \ce{Ca+} \\[-5pt] \ce{Ca+}
\end{array} 
\ce{+ 2Al(OH)3 v} 
\end{align*}

\ce{Ca^2+} 对于土壤胶体还有凝聚作用，有利于形成凝胶，使土壤结构得到改善。

又如，当可溶性肥料，如铵态氮肥或钾肥施入土壤的时侯，土壤溶液里的铵离子和钾离子的浓度增加了。
土壤胶体中原来吸附着的阳离子就被它们所交换。
\begin{align*}
\text{（土壤胶体）}\!\ce{Ca^2+ + 2NH4+ }&\ce{<=>}\text{（土壤胶体）}\! 
\begin{array}{c} \ce{NH4+}\\[-5pt] \ce{NH4+} 
\end{array} 
+ \ce{Ca^2+ } \\
\text{（土壤胶体）}\!\ce{Mg^2+ + 2K+ }&\ce{<=>}\text{（土壤胶体）}\! 
\begin{array}{c} \ce{K+}\\[-5pt] \ce{K+} 
\end{array} 
+ \ce{Mg^2+ } 
\end{align*}

这样，施入肥料里离子态的营养元素就被土壤胶体所吸附，不致因作物没有及时吸收而流失。

当植物根系吸收离子态的营养元素时，使根系周围土壤溶液里营养元素的离子浓度降低，同时因根部放出有机酸和二氧化碳，使土壤溶液里氢离子浓度增加，就把被土壤吸附着的营养元素的离子交换出来，供植物吸收。
\begin{align*}
  \text{（土壤胶体）}\!\ce{NH4+ + H+}&\ce{ <=> }\text{（土壤胶体）}\!\ce{H+ + NH4+ }\\
  \text{（土壤胶体）}\!\ce{K+ + H+ }&\ce{ <=> }\text{（土壤胶体）}\!\ce{H+ + K+} 
\end{align*}

因此，土壤胶体的离子吸附和交换过程对作物营养成分的储蓄和供应起着重要的作用。

\subsection{土壤酸碱性}
农作物的生长对于土壤的酸碱性有一定的要求，有的喜酸性，有的喜碱性，但大多数作物以在中性或弱酸性、弱碱性的土壤中生长为适宜。
适于一般农作物生长的 pH 值在 \numrange{6.5}{7.5} 之间，土壤的酸性或碱性过强都不利于作物的生长。
因此，认识土壤的酸碱性和土壤酸碱性跟农作物生长的关系对农业生产来说是很重要的。

\subsubsection{土壤为什么会显酸性呢？}
植物根部呼吸时放出 \ce{CO2}，跟水反应生成碳酸；土壤有机质的分解，生成多种有机酸，如乙酸、丁酸等；土壤腐殖质所含腐殖酸也是弱酸。
这些物质都使土壤显酸性。
土壤里的铝离子常被土壤胶体吸附，它可以通过离子交换作用进入溶液里,发生水解反应，产生 \ce{H+}，也能使土壤显酸性。
\[ \ce{ Al^3+ + 3H2O <=> Al(OH)3 + 3H+ }\]

酸性土壤可以用施石灰等方法来改良。

\subsubsection{土壤为什么会显碱性呢？}
土壤的碱性主要是由于土壤中 \ce{Na2CO3}、\ce{NaHCO3} 等盐类的水解而产生的。
\begin{gather*}
  \ce{ CO3^2- + H2O <=> HCO3- + OH- }\\
  \ce{ HCO3- + H2O <=> H2CO3 + OH- }
\end{gather*}

碱性土壤可用石膏配合有机肥料、腐殖酸类肥料等来改良。

\subsubsection{土壤的酸碱性跟作物营养的关系}
土壤的酸碱性跟作物营养的关系很大。
如酸性土壤里存在较多的 \ce{Fe^3+} 和 \ce{Al^3+}，当施入可溶性的过磷酸钙等磷肥后，这些离子就跟磷肥起反应生成不溶性磷酸盐，磷肥的肥效就大为降低。
\begin{gather*}
  \ce{ Fe^3+ + PO4^3- = FePO4 v}\\
  \ce{ Al^3+ + PO4^3- = AlPO4 v}
\end{gather*}

农业生产上，在酸性土壤施用可溶性磷肥前，常先施石灰以中和土壤酸性。

土壤的酸碱性会影响土壤里营养成分的有效性。
一般氮的有效性以在 pH 值为 \numrange{6}{8} 之间最高，钾的有效性以在 pH 值为 6 以上最高，磷的有效性以在 pH 值为 \numrange{6.5}{7.5} 时最高。

土壤有机态养分转化为可供作物直接吸收的无机态养分，主要是在微生物参与下进行的（微生物能分泌使有机物分解的酶）。
土壤的酸性或碱性过强，都对微生物生长不利。
例如，施用到土壤里的尿液和尿素化肥，需要在微生物所产生的脲酶的作用下，才会进行水解产生碳酸铵。
中性土壤有利于这一过程的进行，在酸性土壤里，这个反应进行较慢。
\[ \ce{ CO(NH2)2 + 2H2O ->[\text{水解}][\text{脲酶}] (NH4)2CO3 }\]

土壤的酸碱性对植物本身也有影响。
土壤碱性强会使细胞里的原生质溶解，破坏植物组织；酸性强也会使原生质变性，并影响酶的活性，影响植物对养分的吸收。

\subsection{土壤里的氧化—还原反应}
\subsubsection{土壤里的氧化剂和还原剂}
土壤里发生着大量的氧化—还原反应。
许多元素象氧、碳、氮、氢、铁、锰、硫等等都经历着电子转移的变化，许多氧化—还原反应是在微生物存在的条件下进行的，因而更为复杂。
在土壤的氧化—还原反应里，哪些物质属于氧化剂，哪些物质属于还原剂呢？

氧气、氧化铁、二氧化锰、硫酸盐等在土壤里都属于氧化剂。
亚铁盐、锰盐、硫化物、糖类、醛类、羧酸类以及一些易于被微生物分解的有机质在土壤里都属于还原剂。
有机物中，醛类、糖类是还原剂，这早已知道。
为什么羧酸在土壤里是还原剂呢？因为在土壤微生物的作用下，羧酸以及一些易于被微生物分解的有机质会被氧化为二氧化碳。

\subsubsection{土壤里的氧化—还原反应}
当土壤通气条件良好时，土壤团粒的空隙里含有氧气，土壤溶液里也溶解着氧气。
这时，氧化反应顺利进行，把土壤里的 \ce{Fe^2+} 氧化为 \ce{Fe^3+}，\ce{Mn^2+} 氧化为 \ce{MnO2}。
同时，土壤微生物的活动也比较活跃，因为有氧化反应提供足够的能量。
土壤好气性微生物活动旺盛，土壤有机质转变为无机态营养元素的进程加速，养分释放也加快。

相反地，当土壤在渍水条件下或在嫌气性发酵条件和还原剂作用下，\ce{SO4^2-} 还原为 \ce{S^2-}，\ce{Fe^3+} 还原为 \ce{Fe^2+}，\ce{MnO2} 还原为 \ce{Mn^2+}。
这种情况下，嫌气性微生物活动旺盛，不但养分释放慢，而且还积累分解的中间产物，对植物生长往往不利。
如 \ce{Fe^2+} 或 \ce{H2S} 积累到一定浓度时就会毒害植物。

\ce{H2S} 对植物的毒害是影响植物吸收养分和水分，特别是影响吸收钾。
\ce{Fe^2+} 对植物的毒害表现在阻碍对磷和钾的吸收，也影响对氮的吸收。

\subsubsection{土壤的氧化—还原反应对植物营养的关系}
在适宜的氧化—还原反应的条件下，土壤里的有机质受了微生物的作用，不断分解，或成为腐殖质积累在土壤里，或把有机态的氮、磷转变为无机态的氮、磷，供植物直接吸收。
如果氧化条件过于强烈，有机质的矿化过程过于迅速，释放出来的养分因植物来不及吸收而流失，使土壤缺乏有效养分。
如果还原条件过于强烈，有机物的矿化过程缓慢，提供有效养分不足，而且可能造成 \ce{Fe^2+}、\ce{H2S} 等无机毒物和丁酸等有机毒物的积累而影响植物生长。

土壤的氧化—还原状况对植物营养有密切的关系，因此有必要采取各种农业措施以调节土壤的氧化—还原状况。
如旱地的松上，水田的排水、晒田是促进氧化条件的措施；旱地灌水和增施有机肥料，水田的渍水都是促进还原条件的措施。


\subsection{氮、磷、钾在土壤里的转化}
\subsubsection{氮元素在土壤里的转化}
土壤里的氮元素有无机态（或矿质态）的和有机态的。
土壤里的有机质、腐殖质、施入的有机肥料都是有机态氮的来源，其中主要的含氮有机物是氨基酸、蛋白质、核酸等等。
土壤里无机态的氮元素，有的是施入的氮肥，有的是由有机态转化来的，存在形态是 \ce{NH4+}、\ce{NO3-}、\ce{NO2-} 等离子和氮气等等。
植物能够直接利用的是 \ce{NH4+}、\ce{NO3-} 等离子。
豆科植物根瘤里的固氮细菌能够利用游离态的氮。

有机态的氮，如蛋白质等有机物中的氮，绝大部分作物不能直接利用，只有经过微生物的分解作用形成简单的氨基酸或酰胺，再进一步分解为氨或水解成铵盐之后，作物才能吸收利用。
所以，常把有机态氮称为迟效性氮。

目前使用的重要氮肥，除尿素等少数肥料外，绝大多数都属于铵态氮或硝态氮。
尿素（\ce{CO(NH2)2}）是一种酰胺态氮，它在土壤里受微生物的作用转化为铵态氮（\ce{(NH4)2CO3}）。
所以，土壤里无机态氮主要是 \ce{NH4+} 和 \ce{NO3-}。

铵态氮是还原态的氮，氮的化合价为 $-3$；硝态氮是氧化态氮，氮的化合价为 $+5$。

铵态氮在土壤里的亚硝化细菌作用下，氧化为亚硝酸及其盐类。
\[ \ce{ 2NH3 + 3O2 ->[\text{亚硝化细菌}] 2HNO2 + 2H2O }\]

亚硝酸及其盐类在硝化细菌作用下，进一步氧化为硝酸及其盐类。
\[ \ce{ 2HNO2 + O2 ->[\text{硝化细菌}]  2HNO3 }\]

当土壤通气不良或有较多有机质，反硝化细菌就活跃起来，把硝态氮变为氮气等逸出，使氮肥受到损失。
反硝化过程是一个还原过程，因为氮从硝态氮的 $+5$ 价降低到游离态氮的 0 价。

了解氦元素在土壤里的转化情况对研究施用氮肥的方法有重要意义。
例如，铵态氮肥在水稻田里要合理深施，才能很好发挥肥效。
否则铵态氮肥积聚在水田的土壤的表层里，在那里由于灌溉水层中含有一定浓度的溶解氧，为硝化细菌的活动提供了条件，从而把铵态氮氧化为硝态氮。
硝态氮不易被土壤胶体吸附而容易随水流失，即使没有流失而进入水稻土较深处的还原层里，硝态氮将受反硝化细菌作用，变为游离态氮而逸出。
近年来，人们把能抑制硝化细菌的氮肥增效剂跟铵态氮肥一起施入土壤，就能减少铵态氮的转化，减少氮肥损失。

土壤里氮素的增加，主要靠施肥、植物根茎残体的分解、根瘤菌和自生固氮菌从空气中直接吸收固定等途径。
此外，还有小量从降雨和地下水带来的氮素。
在农业生产上，通常采用种植豆科作物，秸杆还田，施用菌肥、有机肥和氮肥增效剂，合理灌溉，改善土壤通气状况，促进有益微生物活动等措施来增加土壤里的氮素。

\subsubsection{磷元素在土壤里的转化}
磷元素在土壤里也以两种形态存在，即无机态磷和有机态磷。
无机态磷主要是钙、镁、铁、铝的磷酸盐。
可以分为：植物不易利用的难溶性磷酸盐，如 \ce{Ca5(PO4)3F}（氟磷灰石）等；植物能利用但必须先行溶解的弱酸溶性磷酸盐，如 \ce{CaHPO4} 等；植物易于利用的水溶性磷酸盐，如 \ce{Ca(H2PO4)2} 等。
有机态磷主要指存在于动植物或微生物体内的含磷有机物，如核蛋白、核酸、磷脂等。
有机态磷绝大部分不能直接被植物吸收利用，而须在微生物作用下，逐步分解为磷酸及其盐类，才能成为植物的养料。
通常我们把有机态磷和难溶性磷酸盐称为迟效性磷，把可溶性磷酸盐以及弱酸溶性磷酸盐称为速效性磷。

可溶性磷酸盐在酸性土壤或碱性土壤里有一部分转变为植物难以利用的磷酸盐，这就是磷的固定作用。
在中性和有机质含量高的土壤里磷素固定作用较弱，磷素的有效性较高。

由于磷的固定作用，使它在土壤里的流动性较氮为小。
因此，磷肥必须施在作物根部附近，才能取得良好的肥效。

\subsubsection{钾元素在土壤里的转化}
钾元素在土壤里的含量要比氮、磷多，但钾元素在土壤里的存在形态主要是复杂的钾的铝硅酸盐，如正长石、水云母等矿物。
这些铝硅酸盐难溶于水，只有在风化过程中，才能逐步转变为能被植物吸收的、可溶的钾的化合物。
所以，钾的铝硅酸盐是迟效性的。
\[ \ce{2KAlSi3O8 + H2O + 2H2CO3 \xlongequal{\quad} Al2(Si2O5)(OH)4 + 2KHCO3 + 4SiO2 } \]

上述风化过程在微生物的作用下，能够加快进行。

土壤中钾的存在形式主要有矿物态钾、土壤胶体吸附的钾和土壤溶液里的钾三种。
其中矿物态钾，占土壤含钾量的 98\% 以上,不易被植物吸收，是迟效性钾。
土壤胶体吸附的钾和土壤溶液里的钾，作物可以直接利用，属于速效性钾。

在土壤风化过程中，在高温多雨的条件下，钾元素比较容易被雨水淋失。
因此，一般说来，我国北方少雨，风化较弱，土壤含钾较丰富；我国南方多雨，风化强烈，土壤含钾较少。
农业生产中，常施用钾肥、草木灰以及有机肥、绿肥等以增加土壤中的钾。

氮、磷、钾三种元素都是作物营养所必需的。
为了实现科学种田，夺取高产，除根据土壤养分状况和作物需肥情况，合理施用氮肥、磷肥外，还要适当配合施用钾肥。
使氮、磷、钾三种营养元素的比率适当。


\chapter{酸、碱和盐的溶解性表（\qty{20}{\celsius}）}
{\noindent\small
\begin{tblr}{colspec={c*{9}{X[c]}},colsep=3pt,rowsep=0pt}
  \diagbox{阳离子}{阴离子} & \ce{OH-} & \ce{NO3-}& \ce{Cl-}& \ce{SO4^2-}& \ce{S^2-}& \ce{SO3^2-}& \ce{CO3^2-}& \ce{SiO3^2-}& \ce{PO4^2-}\\
  \ce{H+}     &        & 溶、挥 & 溶、挥 & 溶 & 溶、挥 & 溶、挥 & 溶、挥 & 微 & 溶 \\
  \ce{NH_4^+} & 溶、挥 & 溶     & 溶     & 溶 & 溶     & 溶     & 溶     & 溶 & 溶 \\
  \ce{K+}     & 溶     & 溶     & 溶     & 溶 & 溶     & 溶     & 溶     & 溶 & 溶 \\
  \ce{Na+}    & 溶     & 溶     & 溶     & 溶 & 溶     & 溶     & 溶     & 溶 & 溶 \\
  \ce{Ba^2+}  & 溶     & 溶     & 溶     & 不 & —      & 不     & 不     & 不 & 不 \\
  \ce{Ca^2+}  & 微     & 溶     & 溶     & 微 & —      & 不     & 不     & 不 & 不 \\
  \ce{Mg^2+}  & 不     & 溶     & 溶     & 溶 & —      & 微     & 微     & 不 & 不 \\
  \ce{Al^3+}  & 不     & 溶     & 溶     & 溶 & —      & —      & —      & 不 & 不 \\
  \ce{Mn^2+}  & 不     & 溶     & 溶     & 溶 & 不     & 不     & 不     & 不 & 不 \\
  \ce{Zn^2+}  & 不     & 溶     & 溶     & 溶 & 不     & 不     & 不     & 不 & 不 \\
  \ce{Cr^3+}  & 不     & 溶     & 溶     & 溶 & —      & —      & —      & 不 & 不 \\
  \ce{Fe^2+}  & 不     & 溶     & 溶     & 溶 & 不     & 不     & 不     & 不 & 不 \\
  \ce{Fe^3+}  & 不     & 溶     & 溶     & 溶 & —      & —      & 不     & 不 & 不 \\
  \ce{Sn^2+}  & 不     & 溶     & 溶     & 溶 & 不     & —      & —      & —  & 不 \\
  \ce{Pb^2+}  & 不     & 溶     & 微     & 不 & 不     & 不     & 不     & 不 & 不 \\
  \ce{Bi^3+}  & 不     & 溶     & —      & 溶 & 不     & 不     & 不     & —  & 不 \\
  \ce{Cu^2+}  & 不     & 溶     & 溶     & 溶 & 不     & 不     & 不     & 不 & 不 \\
  \ce{Hg+}    & —      & 溶     & 不     & 溶 & 不     & 不     & 不     & —  & 不 \\
  \ce{Hg^2+} & —      & 溶     & 溶     & 溶 & 不     & 不     & 不     & —  & 不 \\
  \ce{Ag+}    & —      & 溶     & 不     & 溶 & 不     & 不     & 不     & 不 & 不 \\
\end{tblr}}

{\smallskip\noindent\footnotesize
说明: “溶”表示那种物质可溶于水，“不”表示不溶于水，“微”表示微溶于水，“挥”表示挥发性，“—”表示那种物质不存在或遇到水就分解了。
}